\section*{Введение}
\addcontentsline{toc}{section}{Введение}

\textbf{Актуальность темы.} В современном мире киноиндустрия занимает важное место в культуре и досуге. Практически 
у каждого пользователя интернета есть любимый фильм и несколько люьимых цитат оттуда. Зрители часто запоминают яркие 
цитаты из фильмов, которые становятся частью повседневной речи, мемов и культурного кода. Однако существующие сервисы 
предлагают либо поиск по готовой базе цитат, либо общие сборники без возможности ведения личной коллекции и анализа 
собственных предпочтений. Таким образом, разработка веб-приложения, позволяющего сохранять любимые цитаты, 
систематизировать их по фильмам и получать статистику, является актуальной задачей. 

\textbf{Степень разработанности проблемы.} Существуют сервисы для поиска цитат (PlayPhrase.me, QuoDB), мобильные 
дневники фильмов с функцией заметок (Film Buff) и вики-сборники (Wikiquote). Однако ни один из них не предоставляет 
пользователю возможность вести личный цитатник с аналитикой (топ фильмов, динамика добавления, распределение по жанрам) 
и импортом/экспортом данных. Таким образом, тема является недостаточно проработанной.


\textbf{Объект исследования} — процессы сохранения, систематизации и анализа цитат из фильмов.

\textbf{Предмет исследования} — методы и программные средства реализации веб-приложения для ведения личного цитатника с системой аналитики.

\textbf{Целью данной работы} является разработка веб-приложение, позволяющего пользователям сохранять 
цитаты из фильмов, управлять ими и получать наглядную статистику по своей коллекции.

